% =====================================================================
% Plantilla del artículo del curso (viernes, /paso-11).
% Todas las cifras salen de ../resultados/numeros.tex, que escribe Stata.
% Regla: en este archivo no se teclea ningún número de un resultado.
% =====================================================================
\documentclass[11pt,a4paper]{article}

\usepackage[T1]{fontenc}
\usepackage{lmodern}
\usepackage[spanish,es-noshorthands,es-tabla,es-nodecimaldot]{babel}
\usepackage[margin=2.5cm]{geometry}
\usepackage{amsmath}
\usepackage{booktabs}
\usepackage{graphicx}
\usepackage[authoryear,round]{natbib}
\usepackage[hidelinks]{hyperref}

% Las cifras del texto: las escribe Stata en el /paso-10
\input{../resultados/numeros.tex}

\title{La entrada de genéricos y el precio de los medicamentos de marca:\\
evidencia de un panel de diez principios activos}
\author{[NOMBRE]\\ \small Curso «Investigación en Economía asistida por IA», Universidad de Huelva}
\date{\today}

\begin{document}
\maketitle

\begin{abstract}
\noindent Este trabajo estudia cómo cambia el precio de un medicamento de marca cuando entran competidores genéricos. Con un panel de \nmedicamentos{} medicamentos y \nobs{} observaciones, estimamos una regresión con efectos fijos de medicamento y de año. La entrada de genéricos se asocia con un precio alrededor de un \caidagenerico\,\% más bajo. \textbf{Los datos son ficticios}: el ejercicio sirve para practicar el flujo de trabajo, no para sacar conclusiones.
\end{abstract}

\section{Introducción}

La competencia de los genéricos es uno de los mecanismos principales para contener el gasto farmacéutico. La literatura documenta que el precio de los genéricos cae a medida que aumenta el número de competidores \citep{reiffen2005}, mientras que el precio del medicamento de marca no siempre baja tras la entrada e incluso puede subir \citep{grabowski1992,frank1997}. Este trabajo aporta un ejercicio con un panel en el que la entrada se produce en años distintos para cada medicamento.

\section{Datos}

Usamos un panel ficticio de \nmedicamentos{} medicamentos de marca observados cada año entre 2019 y 2023 (\nobs{} observaciones). Para cada medicamento y año disponemos del precio por envase, el número de envases vendidos, si ya hay genéricos en el mercado y el número de laboratorios que lo comercializan.

\section{Estrategia empírica}

Estimamos
\begin{equation}
\log(\text{precio}_{it}) = \beta_1\, \text{genérico}_{it} + \beta_2\, \text{laboratorios}_{it} + \alpha_i + \gamma_t + \varepsilon_{it},
\end{equation}
donde $\alpha_i$ son efectos fijos de medicamento y $\gamma_t$ efectos fijos de año. Los primeros recogen las diferencias permanentes entre productos; los segundos, los cambios comunes a todos. El diseño no permite una lectura causal: el momento de la entrada puede responder a características del mercado que también afectan al precio.

\section{Resultados}

\begin{table}[htbp]
\centering
\caption{Entrada de genéricos y logaritmo del precio del medicamento de marca}
\label{tab:principal}
\input{../resultados/tablas/tabla1.tex}
\end{table}

\begin{figure}[htbp]
\centering
\includegraphics[width=0.8\textwidth]{../resultados/figuras/figura1.pdf}
\caption{Precio medio (en logaritmos) de los medicamentos que reciben genéricos y de los que no}
\label{fig:precios}
\end{figure}

% [RESULTADOS] Este párrafo lo redacta Claude en el /paso-11,
% usando solo las macros de ../resultados/numeros.tex.

\section{Conclusiones}

En este panel ficticio, la entrada de genéricos va asociada a precios de marca más bajos. Con datos reales, una conclusión así exigiría una estrategia de identificación creíble, comprobar la robustez de los resultados y discutir por qué difieren de los de \citet{grabowski1992}.

\bibliographystyle{apalike}
\bibliography{referencias}

\end{document}
